% Cambios en SN_Payload.tex necesarios para que el manuscrito coincida con paper_revision/
% (respuesta al Comentario 1 del Revisor 1). Cada bloque indica que reemplaza.

%% ---------------------------------------------------------------------------
%% 1. Supuesto 3 y Ec. href: cambiar el signo (hoy dice -(l/g))
%% ---------------------------------------------------------------------------
% En Assumption~\ref{as:smallangle}:  \mathbf{h}_h^\star=(l/g)\,\ddot{\mathbf{x}}_{d,h}
% Reemplazar Ec. (href) y la frase "obtained by filtered numerical differentiation":
\begin{equation}\label{eq:href}
\mathbf{h}_h^\star = \frac{l}{g}\,\ddot{\mathbf{x}}_{d,h},\qquad
\dot{\mathbf{h}}_h^\star = \frac{l}{g}\,\dddot{\mathbf{x}}_{d,h}.
\end{equation}
Both quantities are evaluated analytically from the reference, which is three times
differentiable by Assumption~\ref{as:reference}. With $\mathbf{h}=\mathbf{x}-\mathbf{x}_l$ the
UAV leads the payload in the direction of the reference acceleration.

%% ---------------------------------------------------------------------------
%% 2. Ec. base_reference: agregar el feedforward
%% ---------------------------------------------------------------------------
\begin{equation}\label{eq:base_reference}
\mathbf{v}_{\mathrm{base}} = \dot{\mathbf{x}}_d + \begin{bmatrix}\dot{\mathbf{h}}_h^\star\\ 0\end{bmatrix}
 + K_{\mathrm{traj}}\,\tilde{\mathbf{x}}.
\end{equation}
Since $\mathbf{x}=\mathbf{x}_l+\mathbf{h}$, the UAV velocity that keeps the payload on the reference
with the cable at its quasi-static inclination is $\dot{\mathbf{x}}_d+\dot{\mathbf{h}}^\star$; the
second term supplies the horizontal part of that velocity.

%% ---------------------------------------------------------------------------
%% 3. Modelo del vehiculo (Ec. ModeloLucas): n = 2 canales en el marco de yaw
%% ---------------------------------------------------------------------------
% Reemplazar el parrafo "being n=4 ..." por:
The velocity commands are applied in the yaw frame of the vehicle. Let
$\mathbf{R}_\psi^{h}\in\mathbb{R}^{2\times2}$ be the upper-left block of \eqref{eq:yaw_rotation}. The
horizontal velocity and the horizontal command of the kinematic controller are expressed in that frame,
\begin{equation}\label{eq:yaw_frame}
\mathbf{v} = (\mathbf{R}_\psi^{h})^{\top}\dot{\mathbf{x}}_h,\qquad
\mathbf{v}_d = (\mathbf{R}_\psi^{h})^{\top}\mathbf{v}_{c,h},\qquad \mathbf{v},\mathbf{v}_d\in\mathbb{R}^{2}.
\end{equation}
The dynamic compensator acts on these $n=2$ channels; the vertical and yaw channels receive
$v_{c,z}$ and $\dot\psi_c$ directly. For each channel $i\in\{x,y\}$, model \eqref{eq:ModeloLucas} reads
$\dot v_i=\kappa_{u,i}u_i-\kappa_{v,i}v_i+\delta_i$, with $\kappa_{u,i},\kappa_{v,i}>0$, and $\delta_i$
collects the cable force, the aerodynamic forces, the rotation of the frame and unmodelled dynamics.
It is linear in $\boldsymbol{\theta}_i=[\,1/\kappa_{u,i}\ \ \kappa_{v,i}/\kappa_{u,i}\,]^{\top}\in\mathbb{R}^2$.

%% ---------------------------------------------------------------------------
%% 4. Ley de control y ley de adaptacion (reemplazan AccionControl y ajuste2)
%% ---------------------------------------------------------------------------
\begin{equation}\label{eq:AccionControl}
u_i = \boldsymbol{\phi}_i\hat{\boldsymbol{\theta}}_i - w_i,\quad
\boldsymbol{\phi}_i = [\,a_i\ \ v_i\,],\quad a_i=\dot v_{d,i}-\beta\tilde v_i,\quad
w_i = K_{1,i}\sqrt{|s_i|}\,\mathrm{sign}(s_i)+\nu_i,\quad \dot\nu_i = K_{2,i}\,\mathrm{sign}(s_i),
\end{equation}
with $s_i=\tilde v_i+\beta\int\tilde v_i\,dt$, $\tilde v_i=v_i-v_{d,i}$. Substituting in the model,
\begin{equation}\label{eq:sdot_theta}
\dot s_i = \kappa_{u,i}\boldsymbol{\phi}_i\tilde{\boldsymbol{\theta}}_i - \kappa_{u,i}w_i + \delta_i,
\qquad \tilde{\boldsymbol{\theta}}_i=\hat{\boldsymbol{\theta}}_i-\boldsymbol{\theta}_i .
\end{equation}
On the sliding mode ($s_i=\dot s_i=0$) the integral state equals the equivalent control,
$\nu_i=\boldsymbol{\phi}_i\tilde{\boldsymbol{\theta}}_i+\delta_i/\kappa_{u,i}$. The parameters are updated with
\begin{equation}\label{eq:ajuste_theta}
\dot{\hat{\boldsymbol{\theta}}}_i=\mathrm{Proj}_{\Theta}\!\Big(-\Gamma\,\boldsymbol{\phi}_i^{\top}\,
\mathrm{sat}_{\bar\nu}(\nu_i)-\sigma\big(\hat{\boldsymbol{\theta}}_i-\boldsymbol{\theta}_0\big)\Big)
\quad\text{if } t>t_{\mathrm{on}}\text{ and }|s_i|\le s_{\mathrm{on}},\qquad
\dot{\hat{\boldsymbol{\theta}}}_i=\mathbf{0}\ \text{otherwise},
\end{equation}
where $\Gamma\in\mathbb{R}^{2\times2}$ is positive definite, $\sigma>0$, $\boldsymbol{\theta}_0$ is the
nominal value, $\mathrm{sat}_{\bar\nu}$ is the saturation at $\pm\bar\nu$ and $\Theta$ is a box with
$\theta_{1}\ge\theta_{1,\min}>0$, so that $\hat\kappa_{u,i}>0$. Every signal on the right-hand side is
available in the loop.

% Agregar despues de la ley (o en el Remark de PE):
\begin{remark}[Identifiability]\label{rem:ident}
Law \eqref{eq:ajuste_theta} is a gradient estimator driven by $\nu_i$. It converges to
$\boldsymbol{\theta}_i$ only if $\boldsymbol{\phi}_i$ is persistently exciting and $\delta_i$ has zero
mean. On the circular reference the speed is constant, so $a_i\approx0$ in the yaw frame and
$\theta_{1,i}$ is not excited. The lemniscate, whose speed varies, excites $\theta_{1,x}$. This is why
the mass sweep of Section~\ref{sec:mass} uses the lemniscate.
\end{remark}

%% ---------------------------------------------------------------------------
%% 5. Subseccion nueva al final de la Secc. 2
%% ---------------------------------------------------------------------------
\subsection{Discrete-time implementation}\label{sec:discrete}
The controller runs at the sampling instants of the simulator, nominally $T_c=20$~ms; the period
$\Delta t_k$ is taken from the simulator clock. The estimator
\eqref{eq:h_raw}--\eqref{eq:h_p_filter} is discretised as
$\hat{\mathbf{h}}_p(k)=\hat{\mathbf{h}}_p(k-1)+a_h\big(\mathbf{h}_{p,\mathrm{raw}}(k)-\hat{\mathbf{h}}_p(k-1)\big)$,
$a_h=1-e^{-\Delta t_k/\tau_h}$, and
$\hat{\mathbf{h}}(k+1)=\hat{\mathbf{h}}(k)+\Delta t_k\hat{\mathbf{h}}_p(k)$, with $\hat{\mathbf{h}}(0)=\mathbf{h}(0)$.
With the values of Table~\ref{tab:params} the discrete estimator has damping $\zeta\approx0.44$ and,
with respect to $\dot{\mathbf{h}}$, unit gain and a delay of about 37~ms at the pendulum frequency.
The integrals $\int\tilde v_i\,dt$, $\nu_i$ and $\hat{\boldsymbol{\theta}}_i$ use the forward Euler rule;
the projection onto $\Theta$ is applied after each step. The derivative $\dot v_{d,i}$ is a backward
difference followed by a first-order filter with time constant 50~ms. No other filter is applied.
The command is limited to $|u_i|\le u_{\max}$ and $\nu_i$ is not integrated while the limit is active
in the direction of its update.

%% ---------------------------------------------------------------------------
%% 6. Seccion 3: parrafo de parametros nominales (reemplaza "initialized with arbitrary values...")
%% ---------------------------------------------------------------------------
The parameters of the controller are listed in Table~\ref{tab:params}. The nominal model parameters
follow from the velocity-command interface of the simulator, which has unit static gain:
$\kappa_{u,i}=\kappa_{v,i}=1/\tau$ with $\tau=0.4$~s, i.e., $\boldsymbol{\theta}_0=[0.4,\ 1.0]^{\top}$.
No identification of the simulator was performed. Each case was run five times, because the real-time
co-simulation is not strictly deterministic, and results are reported as mean $\pm$ standard deviation.
% Borrar: "All runs were deterministic, with fixed references and prescribed wind profiles."

%% ---------------------------------------------------------------------------
%% 7. Ec. (44): alpha con el signo de h = x - x_l (Comentario 5, pero afecta las tablas de C1)
%% ---------------------------------------------------------------------------
% \alpha(t)=\mathrm{atan2}\big(\sqrt{h_x^2+h_y^2},\,h_z\big),\qquad \mathbf{v}_{rel}=\dot{\mathbf{x}}_l-\dot{\mathbf{x}}=-\dot{\mathbf{h}}
